\section{Transcendental Desires}

\begin{quotex}
Being endowed with a pure intellect, and controlling oneself with fortitude, rejecting the objects beginning from sound, and eliminating attachment and hatred; one who resorts to solitude, eats sparingly, has speech, body and mind under control, to whom meditation and concentration are ever the highest duty, and who is possessed of dispassion; that person, having discarded egotism, force, pride, desire, anger and superfluous possessions, free from the idea of possession, and serene, is fit for becoming Brahman.
\flright{\textsc{Bhagavad Gita 18:51-53}}
\end{quotex}


Animals are solely motivated by fear, sex, and hunger: that is, the desires for life, reproduction, and food. Animals are constrained to the sensual life. Human beings, on the other hand, are motivated by five transcendental desires which surpass strictly biological needs. These are the desire for:

\begin{itemize}


\item Truth

\item Love

\item Justice

\item Beauty

\item Being
\end{itemize}


\paragraph{Animal Man}


\begin{quotex}
But the sensual or animal man perceiveth not these things that are of the Spirit of God; for it is foolishness to him, and he cannot understand, because it is spiritually examined. But the spiritual man judgeth all things; and he himself is judged of no man.
\flright{\textsc{I Corinthians 2: 14-15}}
\end{quotex}


In our time, it is considered outré to refer to human beings as animals. Yet here we have St Paul asserting the idea that some men live at the level of animals. Moreover, the \emph{Philokalia} is replete with similar claims.

Even the Buddha recognized that there is something beastly deep within man. Three beasts represent the weakness in the heart of creation:

\begin{itemize}


\item \textbf{Ignorance}, represented by the pig


\item \textbf{Malice}, represented by the snake


\item \textbf{Concupiscence}, represented by the rooster

\end{itemize}


Ignorance, malice, and concupiscence are not coincidentally three effects of the Fall. To the extent that these three qualities dominate consciousness, the transcendental desires fade away from awareness.

\paragraph{The Desire for Truth}


Philosophers and scientists, at their best, have a desire to know the Truth. One fairly recent example is \textbf{Brand Blanshard} who, standing on the shoulders of Benedict Spinoza and the British Idealists, wrote a massive study on the \emph{Nature of Thought}. He was able to attain some understanding of the Absolute. Unfortunately, he knew that the Absolute was Truth yet he failed to understand it was also the Good. Thus, he denied the existence of maliciousness. All evil, therefore, he claimed was solely the result of ignorance. All that thinking could not reach the conclusion that the Absolute was also the Good -- or Love, for that matter.

The Maverick Philosopher\footnote{\url{https://maverickphilosopher.typepad.com/maverick\_philosopher/2019/01/could-it-be-like-this.html}}, despite all his intelligence and erudition, after 15 years and 5,000 posts comes to this conclusion:

\begin{quotex}
I find the following scenario exceedingly strange. We die and become nothing and no question gets answered. Could it be like this? It is epistemically possible, possible for all we know. All we know is damned little. But then what would have been the point of the evolution of animals that pose unanswerable questions? No point! Human life would then be like a joke, but a joke without a teller. \textbf{We can't know that the above scenario is true, and we can't know that it is false}. So in the end you must decide what you will believe and how you will live. \textbf{There is no theoretical resolution of the problem}; the resolution must be personal, pragmatic, and existential.
\end{quotex}


If there is no theoretic resolution of the problem, a truly intelligent man would abandon that line of approach and look for another.

\paragraph{The Two Intellectual Centres}


\begin{quotex}
Beware lest any man cheat you by philosophy, and vain deceit; according to the tradition of men, according to the elements of the world, and not according to Christ: For in him dwelleth all the fulness of the Godhead corporeally.
\flright{\textsc{Colossians 2: 8-10}}
\end{quotex}


Apparently, the addiction to rational philosophy is tough to give up. A fortiori, they even presume that God Himself thinks like a philosopher. For example, \textbf{Alvin Plantinga} asserts his in regard to God's omniscience:

\begin{quotex}
If God is omniscient, then He is unlimited in knowledge; He knows every true proposition and believes none that are false.
\end{quotex}


Of course not, God's knowledge is not like that at all. Propositions are statements about reality, so they are once removed; knowledge of propositions is not knowledge of reality. As useful as rational thinking may be, the deepest thinkers, e.g. \textbf{Thomas Aquinas}, regarded it as an inferior form of knowledge. The knowledge of God and the angels is of a higher order, more like \textbf{intuition}, which is direct, unmediated, and nondual. So if rational thought alone cannot provide the solution to life, can humans attain to intuition like the angels. Mystics and saints have indeed achieved such states. Here is an example from \textbf{Saint John of the Cross}.

\begin{quotex}
when a person has finished purifying and voiding himself of all forms and apprehensible images, he will abide in this pure and simple light, and be perfectly transformed into it. This light is never lacking to the soul, but because of creature forms and veils weighing upon and covering it, the light is never infused. If a person will eliminate these impediments and veils, and live in pure nakedness and poverty of spirit ... his soul in its simplicity and purity will then be immediately transformed into simple and pure Wisdom, the Son of God.
\flright{\textsc{John of the Cross}, \emph{The Ascent of Mount Carmel}}
\end{quotex}


John even provides precise instructions: purify and void yourself of all forms and apprehensible images. Despite that, the philosopher\footnote{\url{https://maverickphilosopher.typepad.com/maverick\_philosopher/2010/03/incarnation-a-subjective-approach.html}} is perplexed:

\begin{quotex}
This is admittedly somewhat murky. More needs to be said about the exact sense of `subjective' and `objective.'
\end{quotex}


I don't think so.

\paragraph{The Desire for Love}


Although \textbf{Saint Philomena} was martyred in the year 301 at the age of 13, her remains were not discovered until 1802 in the Catacomb of Priscilla. Nothing was known about her until she revealed the story of her life and death to \textbf{Sister Maria Luisa di Gesu} in 1833. Philomena, the daughter of a Greek king, took a vow as a consecrated virgin at the age of 13.

Having refused the amorous advances of Emperor Diocletian, she was subjected to torments including drowning and being shot with arrows. Miraculous interventions saved her each time until she was eventually beheaded.

The revelation to Sister Luisa demonstrates this principle enunciated by \textbf{Valentin Tomberg}:

\begin{quotex}
It sets in opposition to evocation by necromancy: the ascent towards the deceased effected by the force of love.
\end{quotex}


In other words, it was nothing like a séance that attempts to evoke the spirit of Philomena. Rather, it was a revelation of the forces of purity, faith, and love.

\paragraph{The Desire for Justice}


\textbf{Pope Saint Hyginus} was a Greek who became the bishop of Rome from 138 to 142. He determined the various prerogatives of the clergy and defined the grades of the ecclesiastical hierarchy for the whole Church, East and West. Many of the early Popes came from the Eastern Empire and obviously had no objections to the office of the Pope at that time. Hyginus died a martyr during the reign of Marcus Aurelius.

\paragraph{The Desire for Being}


Our essence is who we are and our personality is what we acquire in life. How, then, do accidents fit into life? Presumably accidents cannot change our essence and we certainly don't deliberately seek them out. \textbf{Aristotle} listed several accidental properties of beings, including in particular space and time.

In a recent discussion on Christian pessimism, someone pondered why so much emphasis was placed on family. In his opinion, most families were ``crud''. Coincidentally, I broached the same topic in a different way when I was visiting grandchildren and nephews. I remarked how those children were fortunate to be born into a loving and intellectually rich environment. So the question arises about the circumstances of one's birth. Are they totally arbitrary, accidental, and random or are they part of our essence?

\textbf{St Augustine} asserted that the predestined are given all they require to attain the beatific vision. That doctrine certainly challenges the notion of invincible ignorance. Moreover, it implies that the circumstances of birth are set: the predestined are born in a place and time where they can hear the gospel, receive the sacraments and so on. Similarly, our families are suitable for our level of being, so there cannot be an injustice in it.

There are five conditions for the manifestation of being: essence, matter, space, time, and life. Essence is transcendent and matter refers to the circumstances of life. We are accustomed to consider that space and time are neutral, featureless containers for being. That is not the case at all. Newton assumed a three-dimensional uniform Euclidean space in which gravitation worked, although he never explained gravity itself. The Theory of Relativity, however, claims that the shape of space-time is warped, which is the explanation for what we experience as gravity.

The Traditional teaching is that there are spiritual beings (angels) of space and time. Thus, various archangels are assigned to geographical areas, families, and nations. Other angels, or demons for that matter, influence the zeitgeist, the spirit of an age. Hence the circumstances of space and time arise not from accidental rearrangements of matter, but have a transcendent source.


\flrightit{Posted on 2019-01-12 by Cologero}

